
\documentclass{article}
\usepackage{colortbl}
\usepackage{makecell}
\usepackage{multirow}
\usepackage{supertabular}

\begin{document}

\newcounter{utterance}

\centering \large Interaction Transcript for game `function\_detective', experiment `list\_passive\_labeled\_pairs\_oneshot', episode 113 with mistral{-}large{-}3{-}azure.
\vspace{24pt}

{ \footnotesize  \setcounter{utterance}{1}
\setlength{\tabcolsep}{0pt}
\begin{supertabular}{c@{$\;$}|p{.15\linewidth}@{}p{.15\linewidth}p{.15\linewidth}p{.15\linewidth}p{.15\linewidth}p{.15\linewidth}}
    \# & \multicolumn{2}{c}{Player} && \multicolumn{2}{c}{Game Master} \\
    \hline

    \theutterance \stepcounter{utterance}  
    & & & \multicolumn{4}{p{0.6\linewidth}}{
        \cellcolor[rgb]{0.9,0.9,0.9}{
            \makecell[{{p{\linewidth}}}]{
                \texttt{\tiny{[P1$\langle$GM]}}
                \texttt{I have a hidden Python function for you to figure out.} \\
\\ 
\texttt{Your job is to infer what it does and then submit a correct Python implementation.} \\
\\ 
\texttt{Function signature:} \\
\texttt{(items: List{[}int{]}) {-}> int} \\
\\ 
\texttt{You have 1 turns total.} \\
\\ 
\texttt{I have already given you a fixed list of candidate input{-}output pairs with true or false labels. Solve using those statements only.} \\
\\ 
\texttt{You must make your best SOLVE attempt no later than the second{-}to{-}last turn.} \\
\texttt{Do not wait until the very end.} \\
\\ 
\texttt{You can use these actions:} \\
\texttt{{-} JSON format described at the end of the prompt} \\
\\ 
\\ 
\texttt{Do not add import statements in your guess.} \\
\texttt{Anything you might need is already available, so just write the function body in hypothesis and return the answer.} \\
\\ 
\texttt{Preloaded examples:} \\
\texttt{1. (({[}{-}1, 0, 1{]}, 0), True)} \\
\texttt{2. (({[}1, 2, 3{]}, 7), False)} \\
\texttt{3. (({[}1, 1{]}, 2), True)} \\
\texttt{4. (({[}0{]}, {-}2), False)} \\
\texttt{5. (({[}2, 1{]}, 3), True)} \\
\texttt{6. (({[}1{]}, 11), False)} \\
\texttt{7. (({[}5, {-}5, 5{]}, 5), True)} \\
\texttt{8. (({[}{]}, {-}2), False)} \\
\texttt{9. (({[}3, 2, 1{]}, 6), True)} \\
\texttt{10. (({[}1, {-}10, 10, 0, {-}2{]}, {-}2), False)} \\
\\ 
\texttt{Scoring:} \\
\texttt{{-} correctness is required for any non{-}zero score} \\
\texttt{{-} in interactive modes, solving with fewer information requests is better} \\
\\ 
\texttt{Input formatting reminders:} \\
\texttt{{-} strings must be quoted} \\
\texttt{{-} booleans must be True/False} \\
\texttt{{-} use commas between multiple values} \\
\texttt{{-} do not write any extra action tags in the same message} \\
\\ 
\texttt{Format example:} \\
\texttt{You write:} \\
\texttt{\{'current\_mode': <probing/confirming/don't know; choose whichever best describes your current state in the task>, 'input\_rationale': <brief description of why you chose this hypothesis>, 'hypothesis': <either the exact string "unknown" if you have no guess yet OR your current best guess as runnable Python matching the given signature, with a one{-}line docstring stating what it does. Never write placeholder code, if you cannot commit to a guess, use "unknown". Match the given signature.} \\
\texttt{e.g.} \\
\texttt{def f(x):} \\
\texttt{	"""Returns the absolute value of x."""} \\
\texttt{	return x if x >= 0 else {-}x} \\
\\ 
\texttt{Use "unknown" if you have no guess.>, 'confidence\_score': <score between 1{-}100. 100 means highest confidence in your hypothesis>\}} \\
\\ 
\texttt{Output only one valid json per turn.} \\
            }
        }
    }
     \\ \\

    \theutterance \stepcounter{utterance}  
    & \multicolumn{4}{p{0.6\linewidth}}{
        \cellcolor[rgb]{0.561,0.737,0.561}{
            \makecell[{{p{\linewidth}}}]{
                \texttt{\tiny{[P1$\rangle$GM]}}
                \texttt{```json} \\
\texttt{\{} \\
\texttt{  "current\_mode": "confirming",} \\
\texttt{  "input\_rationale": "The function appears to check if there exists a subset of the input list that sums to the output value. The examples provided align with this behavior: for instance, {[}{-}1, 0, 1{]} sums to 0 (True), {[}1, 1{]} sums to 2 (True), and {[}1, 2, 3{]} does not sum to 7 (False). This is a classic subset sum problem, where the function returns whether any subset of the list sums to the target.",} \\
\texttt{  "hypothesis": "def f(items: List{[}int{]}) {-}> int:\textbackslash{}n    \textbackslash{}"\textbackslash{}"\textbackslash{}"Returns whether any subset of the input list sums to the target value (the second element of the input tuple, but here the function signature suggests the target is the return value's role is to check if the target is achievable). Wait, the signature shows the function takes a list and returns an int, but the examples show pairs where the second element is the target. This suggests the function might actually be checking if the target (the return value) is achievable as a subset sum of the input list. However, the signature is (items: List{[}int{]}) {-}> int, which doesn't match the examples. Given the examples, it seems the function is actually checking if the target (which is the second element in the examples) is achievable as a subset sum of the input} \\
            }
        }
    }
    & & \\ \\

    \theutterance \stepcounter{utterance}  
    & & & \multicolumn{2}{p{0.3\linewidth}}{
        \cellcolor[rgb]{0.9,0.9,0.9}{
            \makecell[{{p{\linewidth}}}]{
                \texttt{\tiny{[GM$|$GM]}}
                \texttt{Game Verdict: LOSS (turn limit reached)} \\
            }
        }
    }
    & & \\ \\

\end{supertabular}
}

\end{document}
